%% ============================================================
%% 分片 27：原文第 10 节「The Bogomolny Equations And The Space Of Hecke Modifications」（原文 glshort.tex 起于第 8120 行）
%% 由 tools/harvmac2tex.py 从 arXiv 源码（hep-th/0604151v3）自动生成。
%% 只译英文散文；公式、编号、\label、\cite 一律不改。
%% ============================================================

\subsection{Monopole Bubbling} \label{sec:monbub}

In section 第~\ref{sec:spaceheck}~节, we found that the natural compactification
$\bar {{\EUY}}(^L\negthinspace w)$ of the space
${{\EUY}}(^L\negthinspace w)$ of Hecke modifications of some given
type $^L\negthinspace w$ is achieved by including Hecke
modifications of various associated types with smaller weights.

We can now see why this is true from the point of view of the
Bogomolny equations.  Consider a solution of the Bogomolny
equations on $\Bbb{R}^3$ with a specified singularity at the
origin.  It corresponds to an equivariant instanton on
$\Bbb{C}^2\cong \Bbb{R}^4$.  The instanton equations on
$\Bbb{R}^4$ are scale-invariant, so a phenomenon can occur that is
called ``bubbling'' in the mathematical literature: an instanton
can shrink to a point.  In the case of an ${\EUF}$-equivariant
instanton, this point will automatically be the origin in
$\Bbb{C}^2$.

When interpreted in three dimensions using the relation
\cite{ref110}\ between equivariant instantons and solutions of
the Bogomolny equations in three dimensions with point
singularities (which we interpret in terms of 't Hooft operators)
this gives a phenomenon that we might call monopole bubbling.
Monopole bubbling (which was also discovered in unpublished work
by Kronheimer related to \cite{ref110}) is relatively unfamiliar
because it does not occur for smooth solutions of the Bogomolny
equations on $\Bbb{R}^3$ (or other complete Riemannian manifolds)
with the usual boundary conditions. On the contrary, smooth
monopoles on $\Bbb{R}^3$ have a discrete topological
classification  and a characteristic size, which involves the
asymptotic behavior at infinity on $\Bbb{R}^3$.  Monopole bubbling
can occur only in the presence of an 't Hooft operator; it is a
process in which a quantized unit of charge becomes concentrated
at a point (the position of the 't Hooft operator) and disappears.

 We will now very briefly use the ADHM construction \cite{ref125} of instantons on $\Bbb{R}^4\cong \Bbb{C}^2$ to determine
how bubbling of equivariant instantons -- or in other words
monopole bubbling -- affects the weight $^L\negthinspace w$ of an
't Hooft operator for the case of gauge group $G=SU(N)$. (A
similar analysis could be made for other groups for which there is
an ADHM construction, namely the classical groups $SO(N)$ and
$Sp(N)$. But for exceptional groups there is no ADHM
construction.) We will see that bubbling changes the weight in
precisely the expected fashion.

We will follow \cite{ref126}, section 3.3. To describe $SU(N)$ instantons on
$\Bbb{R}^4$, let $U$ be a two-dimensional vector space (which
arises as one of the spin bundles of $\Bbb{R}^4$) on which
${\EUF}$ acts by the matrix % 式 (10.6)
\setcounter{equation}{5}
\begin{equation}\label{eq:normalo}
\left(\begin{matrix}e^{i\theta}
& 0 \\ 0 & e^{-i\theta}\\\end{matrix}\right).
\end{equation}
  For any vector space $V$
with ${\EUF}$ action, we write $\chi(V)$ for the character of
${\EUF}$; for instance,
 $\chi(U)=e^{i\theta}+e^{-i\theta}$. To construct instantons of
instanton number $k$, let ${{\EUH}}$ be a $k$-dimensional vector
space.  And let $E_\infty$ be a vector space of rank $N$ (which
turns out to be the fiber at infinity of the instanton bundle).
The instanton is built from a special kind of complex
% 式 (10.7)
\setcounter{equation}{6}
\begin{equation}\label{eq:onnos}
0\,\underarrow{~~~~}\,{{\EUH}}\underarrow{\,\,\alpha(z)\,\,}{{\EUH}}\otimes
U\oplus E_\infty\underarrow{\,\,\beta(z)\,\,}
{{\EUH}}\,\underarrow{~~~~}\, 0,
\end{equation}
 where $\beta(z)\alpha(z)=0$. For
our limited purposes here, we do not need the details. $\alpha(z)$
and $\beta(z)$ depend linearly on the complex coordinates
$z_1,z_2$ of $\Bbb{C}^2$.  For a smooth instanton, the only
nonzero cohomology group of this complex, for any $z=(z_1,z_2)\in
\Bbb{C}^2$, is $H^1(z)={\rm ker}\,\beta(z)/{\rm im}\,\alpha(z)$.
The fiber at $z$ of the instanton bundle $E$ derived from the ADHM
data is $E_z=H^1(z)$.

To get an ${\EUF}$-equivariant instanton, we pick an action of
${\EUF}$ on $E_\infty$ and ${{\EUH}}$, and choose the complex
\eqref{eq:onnos}\ to be ${\EUF}$-equivariant. The fiber of the instanton
bundle at the origin, which we denote $E_0$, is in particular the
first cohomology group of the corresponding complex:
% 式 (10.8)
\setcounter{equation}{7}
\begin{equation}\label{eq:otonos}
0\longrightarrow{{\EUH}}\underarrow{\,\,\alpha(0)\,\,}{{\EUH}}\otimes
U\oplus E_\infty\underarrow{\,\,\beta(0)\,\,}
{{\EUH}}\longrightarrow 0,
\end{equation}
 For an equivariant complex
$0\underarrow{~~~} V_0\underarrow{\,\,\alpha\,\,}
V_1\underarrow{\,\,\beta\,\,} V_2\underarrow{~~~} 0$ whose only
cohomology is $H^1$, the character of $H^1$ can be expressed as
$\chi(H^1)=-\chi(V_0)+\chi(V_1)-\chi(V_2)$. So in particular
% 式 (10.9)
\setcounter{equation}{8}
\begin{equation}\label{eq:onsno}
\chi(E_0)=\chi(E_\infty)+\chi({{\EUH}})(e^{i\theta}+e^{-i\theta}-2).
\end{equation}

$\chi(E_0)$ is important because, according to Kronheimer's
construction, it determines the weights of the 't Hooft operator
that appears at the origin of $\Bbb{R}^3=\Bbb{C}^2/{\EUF}$ when
the equivariant instanton is interpreted in three dimensions. If
$\chi(E_0)=\sum_{i=1}^N\exp({im_i\theta})$, then  the weights are
$^L\negthinspace w=(m_1,m_2,\dots,m_N)$.

Instanton bubbling occurs when $\alpha$ and $\beta $ are varied in
such a way that, for some decomposition ${{\EUH}}={{\EUH}}_1\oplus
{{\EUH}}_2$, the complex \eqref{eq:onnos}\ decomposes as the direct sum of a
zero complex % 式 (10.10)
\setcounter{equation}{9}
\begin{equation}\label{eq:teto}
0\,\underarrow{~~~~}\,
{{\EUH}}_1\underarrow{\,\,\,0\,\,\,}{{\EUH}}_1\otimes
U\underarrow{\,\,\,0\,\,\,}{{\EUH}}_1\,\underarrow{~~~~}\, 0
\end{equation}
 and
a complex
% 式 (10.11)
\setcounter{equation}{10}
\begin{equation}\label{eq:ponos}
0\,\underarrow{~~~~}\,{{\EUH}}_2\underarrow{\,\,\alpha(z)\,\,}{{\EUH}}_2\otimes
U\oplus E_\infty\underarrow{\,\,\beta(z)\,\,}
{{\EUH}}_2\,\underarrow{~~~~}\, 0
\end{equation}
 with cohomology only in
dimension 1. The ``bubbled'' instanton bundle $E'$, whose
instanton number is reduced by the rank of ${{\EUH}}_1$, has for
its fiber the first cohomology group of this second complex, and
by analogy with \eqref{eq:onsno}, the character of its fiber at the origin
is
$\chi(E'_0)=\chi(E_\infty)+\chi({{\EUH}}_2)(e^{i\theta}+e^{-i\theta}-2)$.
Hence
% 式 (10.12)
\setcounter{equation}{11}
\begin{equation}\label{eq:ponzo}
\chi(E'_0)=\chi(E_0)+\chi({{\EUH}}_1)(2-e^{i\theta}-e^{-i\theta}).
\end{equation}

This formula determines the change in weights under instanton
bubbling.  The weights $m_i$ of $E_0$ and the weights $m'_i$ of
$E'_0$ are related by
% 式 (10.13)
\setcounter{equation}{12}
\begin{equation}\label{eq:yhog}
\sum_{i=1}^N\exp({im'_i\theta})=
\sum_{i=1}^N\exp({im_i\theta})+\chi({{\EUH}}_1)(2-e^{i\theta}-e^{-i\theta}).
\end{equation}

 For example, if ${{\EUH}}_1$ is
one-dimensional with character $\exp(im\theta)$ for some integer
$m$, then the list of weights of $E_0$ must include $m+1$ and
$m-1$ (since $\chi(E'_0)$ is a Fourier series with nonnegative
coefficients). \eqref{eq:yhog}\ implies that precisely two weights change in
this bubbling process, by $(m+1,m-1)\to (m,m)$. This corresponds
to the simplest degeneration of 't Hooft/Hecke operators that we
discussed in section 第~\ref{sec:spaceheck}~节. For another example, suppose
${{\EUH}}_1$ has rank two, with character $\exp({im
\theta})+\exp({im'\theta})$.  If $|m-m'|>1$, we just get a pair of
moves of the kind just described, with weights
$(m+1,m-1,m'+1,m'-1)$ replaced by $(m,m,m',m')$. But if $m'=m+1$,
then \eqref{eq:yhog}\ allows another move with weights $(m+2,m-1)$ replaced
by $(m+1,m)$. This corresponds to another expected degeneration of
't Hooft/Hecke operators. More generally, if ${{\EUH}}_1$ has rank
$s$ and weights $m,m+1,\dots,m+s-1$, then \eqref{eq:yhog}\ leads to a
move\footnote{A complex of the appropriate ADHM form to realize this
move can be achieved with ${{\EUH}}={{\EUH}}_1$, ${{\EUH}}_2=0$
and $E'$ being the trivial bundle with fiber $E_\infty$. With the
characters of $E'_0$ and ${{\EUH}}_1$ given in the text, there is
an essentially unique ${\EUF}$-invariant choice of
$\tau_1,\tau_2,\sigma,$ and $\pi$ that obeys eqn. (3.3.11) of
\cite{ref126}\ and gives a smooth instanton bundle $E$, with instanton
number $s$, that can degenerate to $E'$ by bubbling.} in which
weights $(m+s,m-1)$ of $E_0$ are replaced by weights $(m+s-1,m)$
of $E'_0$. The expected degenerations of 't Hooft/Hecke operators
for $G=SU(N)$ are all compositions of moves of this type.  For a
general $G$, the expected degenerations can be generated by
bubbling of $SU(2)$ instantons embedded in various subgroups of
$G$.

Note in particular that while in the absence of
${\EUF}$-invariance, instantons can be assumed to bubble one at a
time, this is not so once one imposes ${\EUF}$-invariance. There
are irreducible processes that involve the simultaneous bubbling
of an arbitrarily large number of ${\EUF}$-invariant instantons.
See \cite{ref112,ref113} for detailed description of some
Bogomolny moduli spaces on $\Bbb{R}^3$ in which monopole bubbling
phenomena can be seen, in a suitable limit.

\par\medskip{\it Hyper-Kahler Structure At Singularities}

In section 第~\ref{sec:spaceheck}~节, we noted in an example that although the
compactified space $\bar {{\EUY}}(^L\negthinspace w)$ of Hecke
modifications of a given type $^L\negthinspace w$ is not a
hyper-Kahler manifold, its singularities do have a hyper-Kahler
structure.  Now we can explain why.  The singularities come from
monopole bubbling. Monopole bubbling is governed by the moduli
space of ${\EUF}$-equivariant instantons on $\Bbb{C}^2=\Bbb{R}^4$.
But instanton moduli spaces on $\Bbb{R}^4$ have a hyper-Kahler
structure inherited from the hyper-Kahler structure of $\Bbb{R}^4$
(and manifest in the ADHM construction).  Since the group ${\EUF}$
preserves the hyper-Kahler structure of $\Bbb{R}^4$, the space of
${\EUF}$-equivariant instantons is hyper-Kahler (and again, the
ADHM construction makes this manifest).

 We can also  argue in three-dimensional
terms, without using the relation to equivariant instantons in
four dimensions, that monopole bubbling, and hence the
singularities of $\bar {{\EUZ}}({}^L\negthinspace
w_1,\dots,{}^L\negthinspace w_n)$, is governed by a hyper-Kahler
moduli space. Indeed, the moduli space of solutions of the
Bogomolny equations on $\Bbb{R}^3$ is hyper-Kahler
\cite{ref109}.\footnote{This can be shown using the same sort of arguments
we use momentarily to exhibit a Kahler structure for the moduli
space of solutions of the Bogomolny equations on $W=I\times C$.
Those arguments also work if $I$ is replaced by $\Bbb{R}$. If in
addition $C=\Bbb{R}^2$, then $W=\Bbb{R}^3$ can be written as
$\Bbb{R}^2\times \Bbb{R}$ in many different ways, leading to a
hyper-Kahler structure, not just the Kahler structure that we
obtain in section 第~\ref{sec:kahstr}~节. } The spaces of $\bar
{{\EUZ}}({}^L\negthinspace w_1,\dots,{}^L\negthinspace w_n)$ do
not come  from solutions of the Bogomolny equations on
$\Bbb{R}^3$, but on the more complicated three-manifold $W=I
\times C$ with rather particular boundary conditions. So $\bar
{{\EUZ}}({}^L\negthinspace w_1,\dots,{}^L\negthinspace w_n)$ is
not hyper-Kahler. However, monopole bubbling is a local
phenomenon. When a unit of monopole charge shrinks to a point, the
local structure does not depend on the details of the
three-manifold in which this is occurring.  So the structure near
the singularity is hyper-Kahler, just as if one is on $\Bbb{R}^3$.

We give further details presently on the deformation and
resolution of these hyper-Kahler singularities.
